\section{Objetivos}
  \begin{frame}
    \frametitle{Objetivos del tema}
      \begin{itemize}
        \item Aprender a utilizar técnicas de diseño de filtros generales (paso baja, paso alta, paso banda y supresor de banda):

        \begin{itemize}
        \normalsize
        \item Continuos: 
        \begin{itemize}
          \small
          \item Obtener la función de transferencia del filtro
          \item Aproximaciones a los filtros ideales
        \end{itemize}

        \item Discretos:
        \begin{itemize}
          \small
          \item Filtros IIR: a partir de las aproximaciones de filtros continuos
          \item Filtros FIR: a partir de la transformada de Fourier con ventana
        \end{itemize}
      \end{itemize}

      \item Aprender a diseñar filtros específicos/a medida

  \end{itemize}
\end{frame}

\begin{frame}
  \frametitle{Planificación: Semana 1}

  \begin{itemize}
    \item Objetivos de la Teoría:
    \begin{itemize}
      \item Conocer el comportamiento en frecuencia de los filtros continuos generales.
      \item Definir los tipos de filtro ideales básicos \\(paso alto, paso bajo, paso banda, supresor de banda).
      \item Caracterizar el comportamiento de un filtro real (especificaciones de diseño)
    \end{itemize}

    \item Objetivos Prácticos:
    \begin{itemize}
      \item Entender las especificaciones de diseño de un filtro.
      \item Saber especificar el diseño de un filtro en relación a un problema práctico de procesamiento de señal
    \end{itemize}
  \end{itemize}
\end{frame}

\begin{frame}
  \frametitle{Planificación: Semana 2}

  \begin{itemize}
    \item Objetivos Teóricos:
    \begin{itemize}
      \item Entender la diferencia entre los filtros ideales y los reales.  
      \item Conocer diferentes implementaciones de filtros continuos (Butterworth, Tchebyshev I y II, Elíptico).
      \item Entender las características de cada implementación.
    \end{itemize}
    \item Objetivos Prácticos:
    \begin{itemize}
      \item Seleccionar el filtro continuo más adecuado (de los estudiados hasta ahora) para unas especificaciones.
      \item Diseñar e implementar un filtro y utilizarlo para procesar una señal.
    \end{itemize}
  \end{itemize}
\end{frame}

\begin{frame}
  \frametitle{Planificación: Semana 3}

  \begin{itemize}
    \item Objetivos Teóricos:
    \begin{itemize}
      \item Aprender técnicas de diseño de filtros discretos tipo IIR.
      \item Entender las aproximaciones de filtros continuos.
      \item Aplicar transformaciones espectrales para diseñar un filtro a partir de otro ya conocido.
      \item Entender las transformaciones espectrales y saber obtener un filtro a partir de otro ya diseñado.
    \end{itemize}
    \item Objetivos Prácticos:
    \begin{itemize}
      \item Seleccionar una transformación espectral adecuada para convertir un filtro en otro del tipo deseado.
      \item Diseñar un filtro discreto IIR por aproximación de uno continuo.
      \item Filtrar una señal discreta con un filtro IIR.
    \end{itemize}
  \end{itemize}
\end{frame}


\begin{frame}
  \frametitle{Planificación: Semana 4}
  \begin{itemize}
    \item Objetivos Teóricos:
    \begin{itemize}
      \item Aprender técnicas de diseño de filtros discretos tipo FIR.
      \item Entender el efecto del truncamiento (fenómeno de Gibbs) en el diseño de un filtro FIR.
      \item Conocer el uso de diferentes tipos de ventana (hamming, barlett, blackman, ...) para mitigar el efecto del truncamiento.
    \end{itemize}
    \item Objetivos Prácticos:
    \begin{itemize}
      \item Entender el efecto del truncamiento (fenómeno de Gibbs) en el diseño de un filtro FIR.
      \item Conocer el uso de diferentes tipos de ventana (hamming, barlett, blackman, ...) para mitigar el efecto del truncamiento.
      \item Diseñar un filtro FIR dadas unas especificaciones de diseño.
      \item Filtrar una señal discreta con un filtro FIR.
    \end{itemize}
  \end{itemize}
\end{frame}

\begin{frame}
  \frametitle{Planificación: Semana 5}
  \begin{itemize}
    \item Objetivos Teóricos:
    \begin{itemize}
      \item Aprender técnicas avanzadas de diseño de filtros.
      \item Conocer las características y aplicaciones de los filtros: 
      \begin{itemize}
        \item de fase cero y paso todas.
        \item FIR de fase lineal.
        \item de fase mínima y máxima.
        \item combados.
      \end{itemize}
    \end{itemize}
    \item Objetivos Prácticos:
    \begin{itemize}
      \item Diseñar filtros de los tipos estudiados.
    \end{itemize}
  \end{itemize}
\end{frame}

